\chapter{Thí nghiệm và kết quả}
\chaptergoal{Thiết kế đánh giá định lượng và định tính cho bài toán truy hồi phim tương đồng, từ đó rút ra kết luận về chất lượng biểu diễn của từng mô hình.}

\section{Thiết kế thực nghiệm}
\subsection{Mục tiêu}
Thực nghiệm tập trung trả lời ba câu hỏi:
\begin{enumerate}[label=\arabic*.]
    \item Việc thay đổi đơn vị biểu diễn (\textit{overview} vs \textit{movie profile}) cải thiện chất lượng đến mức nào?
    \item Họ mô hình embedding có mang lại cải thiện đáng kể so với mô hình thống kê không?
    \item Những dạng lỗi nào còn tồn tại sau khi đã chuyển sang biểu diễn ngữ nghĩa?
\end{enumerate}

\subsection{Các cấu hình so sánh}
Các cấu hình được tổ chức theo lộ trình tăng dần độ phức tạp:
\begin{itemize}
    \item C1: TF-IDF trên \textit{overview};
    \item C2: TF-IDF trên \textit{movie profile};
    \item C3: Word2Vec aggregation trên \textit{movie profile};
    \item C4: Sentence Transformer trên \textit{movie profile}.
\end{itemize}

Thiết kế này cho phép tách riêng hai nguồn cải thiện: cải thiện do \textit{dữ liệu đầu vào} và cải thiện do \textit{họ mô hình}.

\subsection{Biến kiểm soát}
Để so sánh công bằng, các biến sau được cố định hoặc ghi nhận rõ:
\begin{itemize}
    \item cùng tập dữ liệu \texttt{core\_*} tại thời điểm audit;
    \item cùng chiến lược split query--candidate;
    \item cùng cấu hình Top-$k$ (đánh giá ở $k=5,10$);
    \item cùng quy tắc tạo \textit{movie profile} cho các mô hình dùng profile.
\end{itemize}

\section{Hệ chỉ số đánh giá}
\subsection{Chỉ số truy hồi}
\begin{itemize}
    \item \textbf{Precision@k}: độ chính xác trong Top-$k$;
    \item \textbf{Recall@k}: khả năng bao phủ mục liên quan;
    \item \textbf{nDCG@k}: chất lượng thứ hạng theo vị trí;
    \item \textbf{Cosine score}: mức độ tách cụm biểu diễn trong không gian vector.
\end{itemize}

\subsection{Nguồn nhãn đánh giá}
Trong giai đoạn đầu, đề tài sử dụng weak supervision dựa trên mức giao thể loại (\textit{shared genres}) để tạo tập mục liên quan. Đây không phải nhãn lý tưởng, nhưng đủ ổn định để so sánh tương đối giữa các mô hình khi chưa có bộ nhãn thủ công quy mô lớn.

\section{Kết quả thực nghiệm}
Kết quả định lượng được xuất tự động từ \texttt{training/evaluation/run\_eval.py} về dạng JSON để giảm sai sót nhập tay. Bảng \ref{tab:model_comparison} là khuôn trình bày chuẩn cho báo cáo cuối.

\begin{table}[H]
\centering
\caption{Bảng so sánh mô hình recommendation}
\label{tab:model_comparison}
\begin{tabularx}{\textwidth}{>{\bfseries}lcccc}
\toprule
Mô hình & Precision@5 & Precision@10 & nDCG@10 & Nhận xét \\
\midrule
TF-IDF overview & N/A & N/A & N/A & Baseline tối giản, dùng làm mốc ban đầu \\
TF-IDF profile & N/A & N/A & N/A & Dự kiến cải thiện nhờ mở rộng profile \\
Word2Vec profile & N/A & N/A & N/A & Nhạy với chất lượng token và độ phủ từ vựng \\
Sentence Transformer & N/A & N/A & N/A & Mô hình cải tiến ưu tiên so khớp ngữ nghĩa \\
\bottomrule
\end{tabularx}
\end{table}

\section{Phân tích định tính}
Ngoài số liệu tổng hợp, báo cáo cần đưa các case study dạng truy vấn phim cụ thể, ví dụ:
\begin{itemize}
    \item phim truy vấn thiên về tâm lý nhân vật;
    \item phim truy vấn thiên về nhịp độ hành động;
    \item phim truy vấn có ít review.
\end{itemize}

Với mỗi case, nên so sánh danh sách Top-$k$ giữa baseline và mô hình cải tiến, sau đó diễn giải vì sao mô hình cải tiến chọn được các ứng viên có liên hệ ngữ nghĩa tốt hơn.

\section{Phân tích lỗi}
Phân tích lỗi được chia thành bốn nhóm chính:
\begin{enumerate}[label=\arabic*.]
    \item \textbf{Sparse review}: phim không có hoặc quá ít review khiến profile thiếu tín hiệu cảm nhận;
    \item \textbf{Metadata bias}: phim có overview ngắn, generic hoặc thiên về marketing;
    \item \textbf{Genre overlap noise}: hai phim chung thể loại nhưng khác biệt lớn về tông truyện;
    \item \textbf{Cold-item effect}: phim mới, ít dữ liệu, khó xếp hạng ổn định.
\end{enumerate}

Mỗi nhóm lỗi cần gắn với một hướng xử lý cụ thể, ví dụ tăng cường metadata, bổ sung keyword/cast hoặc dùng reranking dựa trên nhiều tầng tín hiệu.

\section{Giới hạn thực nghiệm}
Một số giới hạn hiện hữu ảnh hưởng trực tiếp đến độ tin cậy của kết quả:
\begin{itemize}
    \item phân bố review lệch mạnh giữa các phim;
    \item thiếu nhiều trường metadata có giá trị cho ranking;
    \item nhãn đánh giá còn ở mức weak supervision;
    \item chưa có vòng lặp đánh giá online từ phản hồi người dùng thực.
\end{itemize}

\section{Kết luận chương}
Kết quả thực nghiệm cần được diễn giải theo tinh thần khoa học: không chỉ nêu mô hình nào tốt hơn, mà phải chỉ ra \textit{vì sao} tốt hơn, \textit{trong điều kiện nào} lợi thế xuất hiện, và \textit{đâu là giới hạn} của kết luận hiện tại. Cách trình bày này tạo nền tảng vững chắc cho các quyết định mở rộng ở chương kết luận.
